% GENERATED FILE. Edit canonical lessons, then run npm run generate:books.
% canonical-source-hash: fnv1a64:d5576eded4a02405
% canonical-lessons: SA-C48-lion, SA-C48-wolf, SA-C48-goat, SA-C48-monkey, SA-C48-goose

\chapter{More Animals}
\label{ch:sa-more-animals}
\hlchaptermodality{48}

\begin{chapteropening}
\textbf{By the end of this chapter:} \emph{I can name a lion, a wolf, a goat, a monkey and a goose in Sanskrit.}

Complete SA-C48-goose and demonstrate the chapter promise.
\end{chapteropening}

\section[siṁhaḥ]{\sk{सिंहः} (siṁhaḥ) — a lion}

\label{lesson:SA-C48-lion}



\begin{quote}
Before the new word: what did \emph{talam} mean?
\end{quote}

\subsection*{You'll want to know: \sk{सिंहः}}
\textbf{\sk{सिंहः}} (\emph{siṁhaḥ}) — \textquotedblleft{}a lion\textquotedblright{}.

\textbf{\sk{सिंहः}} is a lion.

You have been handed this word twice without being told it. \emph{Siṁhapura}, the lion town, is what \emph{Singapore} was worn down from; and \emph{Narasiṁha}, the man-lion, sets it beside \sk{नरः}, the man you can already say. Here is the word on its own at last.

Its own parentage is unsettled — no cousin outside the family has been agreed on. How far it travelled is not in doubt: the language and the people of Sri Lanka are called Sinhala after it. Ends in \textbf{-\sk{अः}}.

\begin{grammarlens}[title={What you've built}]
A word you had already met twice, now yours to say.
\end{grammarlens}

\subsection*{Your turn}
Say these aloud:

\begin{itemize}
\raggedright
  \item \emph{siṁhaḥ}
  \item it once more, slowly
  \item \emph{talam}, then \emph{siṁhaḥ}
\end{itemize}

\begin{tcolorbox}[breakable,colback=teal!4,colframe=teal!35!black,title={Before you move on}]
What does \sk{सिंहः} mean? (\textquotedblleft{}a lion\textquotedblright{}.) Say it once more.
\end{tcolorbox}
\section[vṛkaḥ]{\sk{वृकः} (vṛkaḥ) — a wolf}

\label{lesson:SA-C48-wolf}



\begin{quote}
Before the new word: what did \emph{siṁhaḥ} mean?
\end{quote}

\subsection*{You'll want to know: \sk{वृकः}}
\textbf{\sk{वृकः}} (\emph{vṛkaḥ}) — \textquotedblleft{}a wolf\textquotedblright{}.

\textbf{\sk{वृकः}} is a wolf.

Where \sk{सिंहः} stops at the border, this one runs the whole way: Latin \emph{lupus}, Greek \emph{lykos}, English \emph{wolf}, Russian \emph{volk}. One animal that kept one name from the Ganges to the Atlantic.

The Sanskrit word is used of a tearing thing more generally too, which fits the root sense of tearing. Ends in \textbf{-\sk{अः}}.

\begin{grammarlens}[title={What you've built}]
Two animals, one with cousins abroad and one without.
\end{grammarlens}

\subsection*{Your turn}
Say these aloud:

\begin{itemize}
\raggedright
  \item \emph{vṛkaḥ}
  \item it once more, slowly
  \item \emph{siṁhaḥ}, then \emph{vṛkaḥ}
\end{itemize}

\begin{tcolorbox}[breakable,colback=teal!4,colframe=teal!35!black,title={Before you move on}]
What does \sk{वृकः} mean? (\textquotedblleft{}a wolf\textquotedblright{}.) Say it once more.
\end{tcolorbox}
\section[chāgaḥ]{\sk{छागः} (chāgaḥ) — a goat}

\label{lesson:SA-C48-goat}



\begin{quote}
Before the new word: what did \emph{vṛkaḥ} mean?
\end{quote}

\subsection*{You'll want to know: \sk{छागः}}
\textbf{\sk{छागः}} (\emph{chāgaḥ}) — \textquotedblleft{}a goat\textquotedblright{}.

\textbf{\sk{छागः}} is a goat.

Like \sk{सिंहः} and unlike \sk{वृकः}, it has no agreed cousin outside India: the goat words of Europe and this one do not line up. Two animals out of the same field, and only one of them carried its name across the family.

That is worth holding on to, because it is what makes \sk{वृकः} remarkable rather than ordinary. Ends in \textbf{-\sk{अः}}.

\begin{grammarlens}[title={What you've built}]
Three animals, and two different fates for a name.
\end{grammarlens}

\subsection*{Your turn}
Say these aloud:

\begin{itemize}
\raggedright
  \item \emph{chāgaḥ}
  \item it once more, slowly
  \item \emph{vṛkaḥ}, then \emph{chāgaḥ}
\end{itemize}

\begin{tcolorbox}[breakable,colback=teal!4,colframe=teal!35!black,title={Before you move on}]
What does \sk{छागः} mean? (\textquotedblleft{}a goat\textquotedblright{}.) Say it once more.
\end{tcolorbox}
\section[vānaraḥ]{\sk{वानरः} (vānaraḥ) — a monkey}

\label{lesson:SA-C48-monkey}



\begin{quote}
Before the new word: what did \emph{chāgaḥ} mean?
\end{quote}

\subsection*{You'll want to know: \sk{वानरः}}
\textbf{\sk{वानरः}} (\emph{vānaraḥ}) — \textquotedblleft{}a monkey\textquotedblright{}.

\textbf{\sk{वानरः}} is a monkey.

Two words you own are sitting inside it: \sk{वनम्}, the forest, and \sk{नरः}, the man. Sanskrit's own commentators read \textbf{\sk{वानरः}} as the forest-man, and that reading has held for two thousand years, whatever later scholars have made of it.

So the word costs you one new shape and pays back two old ones. Ends in \textbf{-\sk{अः}}.

\begin{grammarlens}[title={What you've built}]
A word built out of two you had already.
\end{grammarlens}

\subsection*{Your turn}
Say these aloud:

\begin{itemize}
\raggedright
  \item \emph{vānaraḥ}
  \item it once more, slowly
  \item \emph{chāgaḥ}, then \emph{vānaraḥ}
\end{itemize}

\begin{tcolorbox}[breakable,colback=teal!4,colframe=teal!35!black,title={Before you move on}]
What does \sk{वानरः} mean? (\textquotedblleft{}a monkey\textquotedblright{}.) Say it once more.
\end{tcolorbox}
\section[haṁsaḥ]{\sk{हंसः} (haṁsaḥ) — a goose}

\label{lesson:SA-C48-goose}



\begin{quote}
Before the last word of the chapter: what did \emph{chāgaḥ} mean, and what did \emph{vānaraḥ} mean?
\end{quote}

\subsection*{You'll want to know: \sk{हंसः}}
\textbf{\sk{हंसः}} (\emph{haṁsaḥ}) — \textquotedblleft{}a goose\textquotedblright{}.

\textbf{\sk{हंसः}} is a goose, and in poetry a swan.

Another secure cousin, as far-travelled as \sk{वृकः}: Latin \emph{ānser}, Greek \emph{khēn}, German \emph{Gans}, English \emph{goose} and \emph{gander}. The bird held on to its name as faithfully as the wolf did.

In Sanskrit poetry it is the bird that can take milk out of water, so the word came to carry a sense of a mind that sorts things — and \sk{दुग्धम्} is already yours to put beside it. Ends in \textbf{-\sk{अः}}.

\begin{grammarlens}[title={What you've built}]
Five animals: a lion, a wolf, a goat, a monkey, a goose.
\end{grammarlens}

\subsection*{Your turn}
Say these aloud:

\begin{itemize}
\raggedright
  \item \emph{haṁsaḥ}
  \item it once more, slowly
  \item all five in order, then \emph{siṁhaḥ} and \emph{haṁsaḥ} together
\end{itemize}

\begin{tcolorbox}[breakable,colback=teal!4,colframe=teal!35!black,title={Before you move on}]
What does \sk{हंसः} mean? (\textquotedblleft{}a goose\textquotedblright{}.) Say it once more.
\end{tcolorbox}
